% Source-derived chapter 02: Perverse filtrations and vanishing cycles
% Original source: pw-conjecture-gl-n
\chapter{Perverse filtrations and vanishing cycles}
\label{pw-conjecture-gl-n-chapter-02}
\source{pw-conjecture-gl-n}

\begin{remark}[Source section]
\label{pw-conjecture-gl-n-chapter-02-source}
\source{pw-conjecture-gl-n}
\source{pw-conjecture-gl-n}
This chapter reproduces the corresponding section of the retrieved
LaTeX source, including its definitions, statements, examples, and
proofs. It has not yet been translated into Lean.
\notready
\end{remark}

% The source section title was: Perverse filtrations and vanishing cycles
In this section, we introduce the notion of \emph{strong perversity} and show its compatibility with vanishing cycle functors (Proposition \ref{prop1.4}). This plays a crucial role in Section 2 in lifting the $P=W$ conjecture sheaf-theoretically and reduce it to the twisted case.

\subsection{Perverse filtrations}\label{sec1.1}
\source{pw-conjecture-gl-n}
Let $f: X \to Y$ be a proper morphism between irreducible nonsingular quasi-projective varieties (or Deligne--Mumford stacks) with $\mathrm{dim}X = a$ and $\mathrm{dim}Y = b$. Let $r$ be the defect of semismallness of $f$:
\[
r: = \mathrm{dim} X \times_YX - \mathrm{dim}X.
\]
In particular, we have $r = a-b$ when $f$ has equi-dimensional fibers. The perverse filtration 
\[
P_0H^m(X, \BQ) \subset P_1H^m(X, \BQ) \subset \dots \subset H^m(X, \BQ)
\]
is an increasing filtration on the cohomology of $X$ governed by the topology of the morphism $f$; it is defined to be
\[
P_iH^m(X, \BQ) := \mathrm{Im}\left\{ H^{m-a+r}(Y, {^\mathbf{p}\tau_{\leq i}} (Rf_* \BQ_X[a-r])) \to H^m(X, \BQ)\right\}
\]
where $^\mathbf{p}\tau_{\leq * }$ is the perverse truncation functor \cite{BBD}. 

\begin{lem}
\source{pw-conjecture-gl-n}
\notready\label{lem0}
\source{pw-conjecture-gl-n}
If $f$ has equi-dimensional fibers, the perverse filtration on $H^m(X,\BQ)$ terminates at $P_mH^m(X, \BQ)$, \emph{i.e.}
\[
P_mH^m(X, \BQ) = H^m(X, \BQ).
\]
In particular, we have $1 \in P_0H^0(X, \BQ)$.
\end{lem}

\begin{proof}
By definition and the decomposition theorem \cite{BBD}, the dimensions of the graded pieces of the perverse filtration are given by
\begin{equation}\label{llm0}
\source{pw-conjecture-gl-n}
\mathrm{dim} \mathrm{Gr}^P_iH^{m}(X, \BQ) = \mathrm{dim} H^{m-i-(a-r)}\left(Y,{^\mathfrak{p}\CH^i(Rf_* \BQ_X[a-r])}\right).
\end{equation}
Since a perverse sheaf, as a complex of constructible sheaves, is concentrated in degrees $[-b, 0]$, it only has non-trivial cohomology in degrees $\geq -b$. In particular, the right-hand side of (\ref{llm0}) is non-trivial only if 
\[
m-i-(a-r) \geq -b.
\]
Since $r=a-b$, this inequality is equivalent to $m \geq i$.
\end{proof}


A cohomology class $\gamma \in H^l(X, \BQ)$ can be viewed as a morphism $\gamma: \BQ_X \to \BQ_X[l]$, which naturally induces 
\begin{equation}\label{induced}
\source{pw-conjecture-gl-n}
\gamma: Rf_* \BQ_X \to Rf_* \BQ_X [l]
\end{equation}
after pushing forward along $f$. For an integer $c\geq 0$, we say that $\gamma \in H^l(X, \BQ)$ has \emph{strong perversity} $c$ with respect to $f$ if its induced morphism (\ref{induced}) satisfies
\begin{equation}\label{assumption}
\source{pw-conjecture-gl-n}
\gamma \left({^\mathbf{p}\tau_{\leq i} Rf_* \BQ_X}\right) \subset {^\mathbf{p}\tau_{\leq i+(c-l)}} \left(Rf_* \BQ_X  [l] \right), \quad \forall i;
\end{equation}
more precisely, the condition (\ref{assumption}) says that the composition
\[
{^\mathbf{p}\tau_{\leq i} Rf_* \BQ_X} \hookrightarrow Rf_* \BQ_X \xrightarrow{\gamma} Rf_* \BQ_X [l]
\]
factors through ${^\mathbf{p}\tau_{\leq i+(c-l)}} \left(Rf_* \BQ_X  [l] \right) \hookrightarrow Rf_* \BQ_X  [l]$ for any $i$.  Notice that $\gamma$ automatically has strong perversity $l$, so this condition is interesting only when $c < l$. Combining Lemma \ref{lem0} and the following lemma, we see that if $f$ has equi-dimensional fibers and $\gamma$ has strong perversity $c$, then $\gamma \in P_cH^*(X, \BQ)$.

\begin{lem}
\source{pw-conjecture-gl-n}
\notready\label{lem1.1}
\source{pw-conjecture-gl-n}
If $\gamma\in H^l(X, \BQ)$ has strong perversity $c$ with respect to $f$, then taking cup-product with $\gamma$ satisfies
\[
\gamma\cup - : H^m(X, \BQ) \to H^{m+l}(X, \BQ), \quad P_iH^m(X, \BQ) \mapsto P_{i+c}H^{m+l}(X, \BQ).
\]
\end{lem}

\begin{proof}
This follows from taking global cohomology for (\ref{assumption}) and noticing that
\begin{align*}
    H^{m-a+r}\left(Y, {^\mathbf{p}\tau_{\leq i+(c-l)}} \left(Rf_* \BQ_X  [l+a-r] \right)\right) & \\ = H^{(m-a+r)+l}&\left(Y, {^\mathbf{p}\tau_{\leq i+c}} \left(Rf_* \BQ_X  [a-r] \right)\right). \qedhere
    \end{align*}
\end{proof}

In general, the perverse filtration $P_\bullet H^*(X, \BQ)$ may not be multiplicative, \emph{i.e.}, for $\gamma_j \in P_{c_j}H^*(X, \BQ)$ (j=1,2,\dots,s), it may not be true that
\[
\gamma_1 \cup \gamma_2 \cup \cdots \cup \gamma_s \in P_{c_1+\cdots+c_s}H^*(X, \BQ);
\]
see \cite[Exercise 5.6.8]{Park}. In fact, it was proven in \cite[Theorem 0.6]{dCMS} that Conjecture \ref{conj} is equivalent to the multiplicativity of the perverse filtration associated with the Hitchin system. The following easy observation illustrates the advantage of considering strong perversity in view of the multiplicativity issue.

\begin{lem}
\source{pw-conjecture-gl-n}
\notready\label{lem1.2}
\source{pw-conjecture-gl-n}
If the class $\gamma_j \in H^{l_j}(X, \BQ)$ ($j=1,2,\dots, s$) has strong perversity $c_j$ with respect to $f$, then the cup product 
\[
\gamma_1 \cup \gamma_2 \cup \cdots \cup \gamma_s \in H^{l_1+\cdots+l_s}(X, \BQ)
\]
has strong perversity $\sum_jc_j$.
\end{lem}


\subsection{Vanishing cycles}\label{Sec1.2}
\source{pw-conjecture-gl-n}

Throughout section \ref{Sec1.2}, we let $g: X \to \BA^1$ be a morphism such that $X$ is nonsingular and irreducible with $X_0 = g^{-1}(0)$ the closed fiber over $0 \in \BA^1$. We consider the vanishing cycle functor 
\[
\varphi_g: D^b_c(X) \rightarrow D^b_c(X_0)
\]
which preserves the perverse $t$-structures,
\[
\varphi_g: \mathrm{Perv}(X) \rightarrow \mathrm{Perv}(X_0).
\]
Here $\mathrm{Perv}(-)$ stands for the abelian category of perverse sheaves. We denote by 
\[
\varphi_g : = \varphi_g(\mathrm{IC}_X) = \varphi_g(\BQ_X[\mathrm{dim}X]) \in \mathrm{Perv}(X_0) \]
the perverse sheaf of vanishing cycles. We use $X' \subset X$ to denote the support of the vanishing cycle complex $\varphi_g$ so that $\varphi_g \in \mathrm{Perv}(X')$.

Recall that for any bounded constructible object $\CK \in D^b_c(X)$, a class $\gamma \in H^l(X, \BQ)$ induces a morphism
\[
\gamma: \CK \to \CK [l] 
\]
via taking the tensor product with $\gamma: \BQ_X \to \BQ_X[l]$. The following lemma shows the compatibility between the vanishing cycle functor and restriction of cohomology classes.

\begin{lem}
\source{pw-conjecture-gl-n}
\notready\label{lem1.3}
\source{pw-conjecture-gl-n}
With the same notation as above, let $i: X'\hookrightarrow X$ be the closed embedding. The morphism \[
i^*\gamma: \varphi_g \rightarrow \varphi_g[l] \in D^b_c(X')
\]
induced by the class $i^*\gamma$ (applied to the object $\CK = \varphi_g$) coincides with 
\[
\varphi_g(\gamma): \varphi_g \to \varphi_g[l] \in D_c^b(X')
\]
obtained by applying the functor $\varphi_g$ to $\gamma: \BQ_X \to \BQ_X[l]$.
\end{lem}

\begin{proof}
Let $\iota: X_0 \hookrightarrow X$ be the closed embedding of the closed fiber over $0$. By \cite[Definition 8.6.2]{KS}, the vanishing cycle functor can be written as
\begin{equation}\label{vanishing}
\source{pw-conjecture-gl-n}
\varphi_g(-) = \iota^* R\CH{om}(\CC, -) \in D^b_c(X_0)
\end{equation}
with $\CC$ a fixed complex of sheaves. The morphism obtained
by applying $R\CH{om}(\CC, -)$ to $\gamma: \BQ_X \to \BQ_X[l]$ 
is equivalent to the morphism induced by applying $\gamma$ to $R\CH{om}(\CC, \BQ_X)$.
Similarly, the functor $\iota^*$ sends the morphism induced by $\gamma$
to the morphism induced by $\iota^*\gamma$.

%by applying (\ref{vanishing}) to $\gamma: \BQ_X \to \BQ_X[l]$ is equivalent to the morphism applied $\iota^*\gamma$ to the object $\iota^* R\CH{om}(\CK, \BQ_X)$:
%\[
%\iota^*\gamma: \iota^* R\CH{om}(\CK, \BQ_X) \to \iota^* R\CH{om}(\CK, %\BQ_X)[l] \in \mathrm{Perv}(X_0).
%\]
Therefore, if we denote by $\iota': X' \hookrightarrow X_0$ the closed embedding and view $\varphi_g$ as a perverse sheaf on $X'$, we have an equivalent morphism
\begin{equation}\label{321} 
\source{pw-conjecture-gl-n}
\varphi_g(\gamma) = \iota^*\gamma:  \iota'_*\varphi_g \to \iota'_*\varphi_g[l] \in D^b_c(X_0). 
\end{equation}


%We note that Betti moduliy definition, the vanishing cycle complex $\varphi_g$, as an oBetti moduliject on $X_0$, is of the form
%\[
%\varphi_g = \iota^* \CF
%\]
%with $\CF$ a constructiBetti modulile oBetti moduliject on $X$. The cohomology class $\gamma \in H^l(X, \BQ)$ then induces 
%\[
%\gamma: \CF \to \CF [l] \in D^b_c(X),
%\]
%whose restriction to $X_0$ induces the vanishing cycle morphism $\varphi_g(\gamma)$ on $X_0$:
%\[
%\varphi_g(\gamma) = \iota^*\gamma:  \iota^* \CF \to \iota^*\CF[l] \in D^b_c(X_0). 
%\]
%Finally, as the object $\iota^*\CF= \varphi_g$ is supporte on $X'$, for the closed embedding $\iota': X' \hookrightarrow X_0$ we have
%\[
%\iota^*\gamma: \iota'_* \varphi_g \to \iota'_* \varphi_g[l] \in D^b_c(X_0).
%\]
Finally, after applying $\iota'^*: D^b_c(X_0) \rightarrow D^b_c(X')$ to (\ref{321}) and noticing $\iota'^*\iota'_*= \mathrm{id}$, we obtain that the class $i^*\gamma = \iota'^*\iota^*\gamma$ induces
\[
\varphi_g(\gamma): \varphi_g \to \varphi_g[l] \in D_c^b(X').
\]
This completes the proof.
\end{proof}


\begin{prop}
\source{pw-conjecture-gl-n}
\notready\label{prop1.4}
\source{pw-conjecture-gl-n}
Let $g: X \to \BA^1$ and $X'$ be as above. Assume that $X'$ is nonsingular and
\begin{equation}\label{4.1_0}
\source{pw-conjecture-gl-n}
\varphi_g \simeq \mathrm{IC}_{X'} = \BQ_{X'}[\mathrm{dim}X'] \in \mathrm{Perv}(X').
\end{equation}
Assume further that we have the commutative diagram
\begin{equation*}
\begin{tikzcd}
X' \arrow[r, "i"] \arrow[d, "f'"]
& X \arrow[d, "f"] \\
Y' \arrow[r]
& Y
\end{tikzcd}
\end{equation*}
such that $f$ is proper and $g =   f\circ \nu$ with $\nu: Y \to \BA^1$. Then if a class $\gamma \in H^l(X, \BQ)$ has strong perversity $c$ with respect to $f$, its restriction $i^*\gamma \in H^l(X', \BQ)$ has strong perversity $c$ with respect to $f'$. 
\end{prop}


\begin{proof}
By definition, the morphism 
\begin{equation}\label{1.4_1}
\source{pw-conjecture-gl-n}
\gamma: Rf_* \BQ_X \to Rf_* \BQ_X[l] \in D^b_c(Y)
\end{equation}
induced by $\gamma$ satisfies
\begin{equation}\label{1.4_2}
\source{pw-conjecture-gl-n}
\gamma \left({^\mathbf{p}\tau_{\leq i} Rf_* \BQ_X}\right) \subset {^\mathbf{p}\tau_{\leq i+(c-l)}} (Rf_* \BQ_X[l]), \quad \forall i.
\end{equation}
Now we apply the vanishing cycle functor $\varphi_\nu$ to (\ref{1.4_1}). On one hand, we have the base change $Rf'_* \circ \varphi_g \simeq \varphi_\nu\circ Rf_*$ and the fact that the vanishing cycle functor preserves the perverse $t$-structures. Therefore (\ref{1.4_2}) implies  that the morphism
\[
\varphi_g(\gamma): Rf'_* \varphi_g \to Rf'_* \varphi_g[l] \in D^b_c(Y')
\]
satisfies
\[
\varphi_g(\gamma) \left({^\mathbf{p}\tau_{\leq i}} Rf'_* \varphi_g\right) \subset {^\mathbf{p}\tau_{\leq i+(c-l)}} (Rf'_* \varphi_g[l]), \quad \forall i.
\]
On the other hand, using the isomorphism (\ref{4.1_0}) and Lemma \ref{lem1.3}, the above equation means precisely that the morphism $i^*\gamma : \BQ_{X'} \to \BQ_{X'}[l]$ satisfies
\begin{equation*}
(i^*\gamma) \left({^\mathbf{p}\tau_{\leq i} Rf'_* \BQ_{X'}}\right) \subset {^\mathbf{p}\tau_{\leq i+(c-l)}} (Rf'_* \BQ_{X'}[l]), \quad \forall i;
\end{equation*}
that is, the class $i^*\gamma$ has strong perversity $c$ with respect to $f'$.
\end{proof}
